%% ma: L12-B18-0021
%% lop: 12
%% chuong: 6
%% bai: 18
%% dang: 12.18.9
%% loai: TN4
%% muc: TH
%% dap_an: "C"
%% nguon: "(NBV)-10 ĐỀ KIỂM TRA HK2-TOÁN 12-DUNG CHUNG 3 SÁCH.pdf (H04, Đề số 4), Phần 1 Câu 9"
%% kien_thuc: "Bài 18 (công thức nhân xác suất)"
%% trang_thai: cho-duyet
%% ly_do: "Dạng mới đề xuất 12.18.9 (Công thức nhân xác suất P(AB) = P(A)P(B|A) trong bài toán thời tiết), chờ giáo viên duyệt."
%% ghi_chu: "Đáp án gốc C đúng. Bỏ ảnh minh họa trời nắng (ảnh trang trí)."
\begin{ex}
Xác suất để ngày thứ Hai trời nắng là $0{,}6$. Nếu ngày thứ Hai trời nắng thì xác suất để ngày thứ Ba trời nắng là $0{,}7$. Tính xác suất để cả hai ngày thứ Hai và thứ Ba trời đều có nắng.
\choice{$\dfrac67$}{$\dfrac57$}{\True $0{,}42$}{$0{,}58$}
\loigiai{Gọi $A$: ``thứ Hai nắng'', $B$: ``thứ Ba nắng''. $P(A)=0{,}6$, $P(B\mid A)=0{,}7$. Khi đó $P(AB)=P(A)\cdot P(B\mid A)=0{,}6\cdot0{,}7=0{,}42$. Chọn C.}
\end{ex}
